\begin{abstract}
Scarce labels remain a practical bottleneck in geospatial image analysis, yet many remote sensing scene classification studies still rely on protocol-specific settings that are hard to reproduce or compare fairly. Rather than introducing another remote-sensing-specific architecture, we present a reproducible adaptive-transfer workflow for scarce-label scene classification built on \textbf{DINOv2-Base}, fixed class-balanced split manifests, and scripted multi-seed evaluation. The model freezes most pretrained parameters, selectively unfreezes only the last transformer block and final normalization, inserts a residual feature adapter, and applies momentum-updated feature-center regularization. Experiments on AID and NWPU-RESISC45 follow a fixed protocol with 20 training images and 10 validation images per class, and all main comparisons are reported as three-seed averages regenerated from the same experiment outputs. Against a fair ConvNeXt-Tiny baseline under identical manifests, the workflow improves mean overall accuracy from 0.9179 to 0.9413 on AID and yields a smaller but still positive increase from 0.8721 to 0.8790 on NWPU-RESISC45. Same-backbone ablations show that selective finetuning provides most of the mean gain, while the adapter mainly reduces seed variance and the center regularizer improves performance consistency. Although the full backbone is not lightweight in total inference cost, only about 7.5 million parameters are trainable and the final setup remains feasible on an 8 GB GPU. The resulting contribution is a practical geocomputing workflow centered on public data, fixed manifests, and rebuildable evidence rather than on architectural novelty alone.
\end{abstract}
